% Generated draft: requires mathematical and source review.
\documentclass[11pt]{article}
\usepackage[margin=25mm]{geometry}
\usepackage{fontspec}
\setmainfont{Latin Modern Roman}
\usepackage{amsmath,amssymb,mathtools,unicode-math}
\setmathfont{Latin Modern Math}
\usepackage{enumitem,needspace,xurl,hyperref}
\setlength{\parindent}{0pt}
\setlength{\parskip}{5pt}
\setlength{\emergencystretch}{3em}
\newcommand{\N}{\mathbb{N}}
\newcommand{\Z}{\mathbb{Z}}
\newcommand{\Q}{\mathbb{Q}}
\newcommand{\R}{\mathbb{R}}
\newcommand{\C}{\mathbb{C}}
\newcommand{\F}{\mathbb{F}}
\newcommand{\A}{\mathbb{A}}
\newcommand{\PP}{\mathbb{P}}
\newcommand{\E}{\mathbb{E}}
\newcommand{\eps}{\varepsilon}
\newcommand{\dd}{\,\mathrm{d}}
\begin{document}
\section*{UnsolvedMath ID: 20001937}
\textbf{Source problem code:} catalogue entry
\subsection{Definitions and mathematical statement}
The source asks to describe all affinely equivalent Finsler metrics. Its supplied verification defines affine equivalence by preserving geodesics with their constant-speed affine parametrizations. Thus the canonical geodesic sprays agree; sharing only unparametrized paths is a different, projective question.

Use smooth strongly convex Finsler metrics on the slit tangent bundle \(TM\setminus0\), allowing nonreversibility unless separately imposed. For \(L=F^2/2\), positivity, positive two-homogeneity and positive-definite vertical Hessian are required. Fix a spray
\[
 S=y^i\frac{\partial}{\partial x^i}-2G^i(x,y)\frac{\partial}{\partial y^i}.
\]
The aid reports an exact energy-level characterization: the affine-equivalent metrics with this spray correspond to the solutions of
\[
 \mathcal E_i^S(L):=
 S\left(\frac{\partial L}{\partial y^i}\right)
 -\frac{\partial L}{\partial x^i}=0
\]
that satisfy the stated positivity, homogeneity and strong-convexity conditions. The equations are linear in \(L\), and the admissible energies form a convex cone. This is a PDE description for arbitrary sprays, not a finite explicit classification of every nonlinear spray.

If \(S\) is the quadratic spray of a torsion-free affine connection on a connected manifold, the aid reports classification by smooth strongly convex Minkowski norms on one tangent space invariant under the full holonomy group. Parallel transport extends such a norm, and full holonomy, including disconnected holonomy, is required for global path independence. If that holonomy preserves a Riemannian norm and acts transitively on its unit sphere, every invariant Finsler norm is a constant positive multiple of that norm.

For an explicit family, suppose a Riemannian tangent bundle has a global orthogonal parallel splitting \(TM=E_1\oplus E_2\) with both summands nonzero. For \(v=v_1+v_2\ne0\), put \(A=g(v_1,v_1)\), \(B=g(v_2,v_2)\). The reported family is
\[
 F_\epsilon(v)^2=A+B+\epsilon\frac{AB}{A+B},\qquad
 |\epsilon|\leq\tfrac13,
\]
with \(F_\epsilon(0)=0\). The aid states that it has the same spray as \(g\), is non-Riemannian when \(\epsilon\ne0\), and its fundamental tensor satisfies
\[
 g^{F_\epsilon}_{x,v}(w,w)
 =\tfrac12D_v^2(F_\epsilon^2)(w,w)
 \geq\frac{49}{162}\,g_x(w,w)
 \quad(v\ne0,\ w\in T_xM).
\]
The zero parameter returns the Riemannian norm. Smoothness here is on the slit bundle; no quadratic smooth extension at the zero section is asserted.

The record also includes a transition to unparameterized Riemannian reconstruction. It cites [Matveev2012a] and [Eas-Mat2008] for a projective/path-structure method using nondegenerate parallel tensor-bundle sections, holonomy preservation and curvature compatibility. The exact bundle, connection and metrizability hypotheses are not supplied, so that context is not substituted for the Finsler affine classification. The final line introduces a further open problem not included in this record.
\subsection{Short English statement}
Affine equivalence preserves the geodesic spray, including parametrization. The aid describes its Finsler energy solutions by linear Euler--Lagrange equations with positivity and convexity constraints, and gives a full-holonomy description for quadratic sprays. It also reports an explicit non-Riemannian family on a parallel Riemannian product, while keeping the adjacent projective reconstruction discussion distinct.
\subsection{Sources}
\begin{itemize}
\item Local source record: catalogue entry, catalogue entry; canonical \texttt{aim-geometry-notes.json}, \texttt{notes[274]}, source item 12.0, Problem 12.0.4. Title: ``Spray and holonomy classifications for affinely equivalent Finsler metrics.''
\item Burns--Matveev's AIM list \emph{Open problems and questions about geodesics}, workshop ``Geodesics,'' supplied reference \url{https://aimath.org/pastworkshops/geodesicsproblems.pdf}. The verification supplies the preceding affine-parametrization definition. This draft uses the supplied text without accessing a PDF; the page header is an extraction artifact.
\item Original transition references [Matveev2012a] and [Eas-Mat2008] for unparameterized Riemannian metrizability. No full bibliographic entries or precise tensor-bundle construction are supplied.
\item Research attempt 1 supplies the spray/holonomy descriptions and quantitative parallel-splitting family. Registry dated 14 August 2026: ``partially\_solved,'' ``PARTIAL-PROGRESS,'' statement status ``exact''; attempt partial, research-run status open, rigor ``proved,'' novelty confidence low. The machine-generated, non-peer-reviewed aid requires independent verification before citation.
\end{itemize}
\end{document}
